\documentclass[11pt, letterpaper]{article}

\usepackage[margin=2.5cm]{geometry}
\usepackage[spanish]{babel}
\usepackage{xcolor}
\usepackage{amsmath,amssymb}
\usepackage{graphicx}
\usepackage{booktabs}
\usepackage{fancyhdr}
\usepackage{tcolorbox}
\usepackage{hyperref}

\definecolor{natexorange}{RGB}{255, 107, 53}
\definecolor{darkslate}{RGB}{15, 23, 42}

\pagestyle{fancy}
\fancyhf{}
\fancyhead[L]{\small \color{gray} Informe de Laboratorio -- Asignatura}
\fancyhead[R]{\small \color{gray} NaTex Studio}
\fancyfoot[C]{\thepage}

\begin{document}

% PORTADA FORMAL UNIVERSITARIA
\begin{titlepage}
  \centering
  {\scshape\LARGE Universidad Nacional de Ingeniería \par}
  \vspace{0.5cm}
  {\scshape\Large Facultad de Ciencias Físicas y Matemáticas \par}
  {\scshape\large Departamento de Ingeniería Eléctrica \par}
  \vspace{3cm}

  \begin{tcolorbox}[colback=darkslate, colframe=natexorange, arc=3mm, boxrule=2pt]
    \centering
    \vspace{0.3cm}
    {\Huge \bfseries \color{white} INFORME DE LABORATORIO N$^\circ$ 1 \par}
    \vspace{0.4cm}
    {\Large \color{natexorange} Análisis de Circuitos RLC y Respuesta en Frecuencia \par}
    \vspace{0.3cm}
  \end{tcolorbox}

  \vspace{3cm}

  \begin{flushleft}
    \large
    \textbf{Integrantes:} \\
    $\bullet$ Juan Pérez González (RUT: 20.123.456-7) \\
    $\bullet$ María Rodríguez Silva (RUT: 19.876.543-2) \\[0.4cm]
    \textbf{Profesor Guía:} Dr. Carlos M. Benítez \\
    \textbf{Ayudante:} Ing. Roberto Soto \\
    \textbf{Fecha de Entrega:} \today
  \end{flushleft}

  \vfill
  {\large Santiago, Chile \par}
\end{titlepage}

\tableofcontents
\newpage

\section{Resumen Ejecutivo}
En esta experiencia práctica se estudió el comportamiento en régimen transitorio y permanente de un circuito RLC serie. Se midió la frecuencia de resonancia natural y el factor de calidad $Q$, comparando los resultados experimentales con las predicciones del modelo teórico. Los resultados arrojaron un error porcentual inferior al $2.8\%$.

\section{Introducción y Objetivos}
\subsection{Objetivo General}
Caracterizar la respuesta frecuencial de un circuito RLC serie excitado por una señal senoidal.

\subsection{Objetivos Específicos}
\begin{enumerate}
  \item Determinar experimentalmente la frecuencia de resonancia $\omega_0$.
  \item Calcular el ancho de banda a $-3\text{ dB}$.
  \item Analizar la influencia de la resistencia en el amortiguamiento del sistema.
\end{enumerate}

\section{Marco Teórico}
La ecuación diferencial que rige la corriente $i(t)$ en el circuito serie es:
\begin{equation}
  L \frac{d^2 i(t)}{dt^2} + R \frac{di(t)}{dt} + \frac{1}{C} i(t) = \frac{dv_s(t)}{dt}
\end{equation}

La frecuencia de resonancia teórica y el factor de calidad vienen dados por:
\begin{equation}
  \omega_0 = \frac{1}{\sqrt{LC}}, \qquad Q = \frac{\omega_0 L}{R} = \frac{1}{R} \sqrt{\frac{L}{C}}
\end{equation}

\section{Procedimiento Experimental y Datos}
Se utilizó un osciloscopio digital, un generador de funciones, un inductor de $L = 10\text{ mH}$, un capacitor de $C = 100\text{ nF}$ y una resistencia de $R = 100\,\Omega$.

\begin{table}[h]
  \centering
  \caption{Mediciones de voltaje y desfase en función de la frecuencia.}
  \vspace{0.2cm}
  \begin{tabular}{ccccc}
    \toprule
    \textbf{Frecuencia (kHz)} & \textbf{$V_{\text{in}}$ (V)} & \textbf{$V_{\text{out}}$ (V)} & \textbf{Ganancia (dB)} & \textbf{Desfase ($^\circ$)} \\
    \midrule
    1.0 & 5.00 & 1.25 & -12.04 & +75 \\
    3.0 & 5.00 & 3.48 & -3.15 & +42 \\
    \textbf{5.03 ($\approx f_0$)} & \textbf{5.00} & \textbf{4.92} & \textbf{-0.14} & \textbf{0} \\
    7.0 & 5.00 & 3.20 & -3.87 & -48 \\
    10.0 & 5.00 & 1.10 & -13.15 & -78 \\
    \bottomrule
  \end{tabular}
\end{table}

\section{Conclusiones}
Se comprobó satisfactoriamente la teoría de resonancia en circuitos RLC. El factor de calidad obtenido experimentalmente fue de $Q_{\text{exp}} = 3.12$, concordando estrechamente con el valor analítico esperado de $Q_{\text{teo}} = 3.16$.

\end{document}
